\documentclass[../principal.tex]{subfiles}
\graphicspath{{\subfix{../imagenes/}}}

\begin{document}

\ifSubfilesClassLoaded{
    \part{Ejemplos}
}{}

\chapter{Cuatro cuadradas}

En este capítulo abordaremos el diseño de un sintetizador modular de cuatro pasos de onda cuadrada, incluyendo una cadena autoncontenida, desde reloj hasta parlante.

A grandes ragos tendremos los siguientes módulos:

\begin{enumerate}
    \item Reloj
    \item Secuenciador
    \item Osciladores
    \item Amplificador
    \item Parlante
\end{enumerate}

Como diagrama de bloques proponemos el siguiente:

\begin{tikzpicture}[node distance=2cm]

\node (reloj) [process] {Reloj};
\node (secuenciador) [process, below of=reloj] {Secuenciador};
\node (osciladores) [process, below of=secuenciador] {Osciladores};
\node (amplificador) [process, below of=osciladores] {Amplificador};
\node (parlante) [process, below of=amplificador] {Parlante};



% \node (pro2b) [process, right of=dec1, xshift=2cm] {Process 2b};
% \node (out1) [io, below of=pro2a] {Output};
% \node (stop) [startstop, below of=out1] {Stop};

\draw[arrow, dashed] (reloj) -- node[left]{tierra} (secuenciador);
\draw[arrow] (reloj) -- node[right]{señal de reloj} (secuenciador);
\draw [arrow] (secuenciador) -- (osciladores);
\draw [arrow] (osciladores) -- (amplificador);
\draw [arrow] (amplificador) -- (parlante);

\end{tikzpicture}

\section{Reloj}

Para el módulo de reloj proponemos el uso del chip 555 en modo astable, explicado en \ref{sec:chip555}.


\begin{circuitikz}
    \draw (0,0) node[draw, minimum height=2.0cm, minimum width=1.0cm] (box) {Reloj};
    \draw (box.north) -- ++(0.0, 0.2) node[anchor=south] {$V_{CC}$};
    \draw (box.south) -- ++(0.0, -0.2) node[anchor=north] {Tierra};
    \draw (box.east) |- ++(0.4,0.2) node[anchor=west] {Señal de reloj};
    \draw (box.east) |- ++(0.4,-0.2) node[anchor=west] {Tierra};
\end{circuitikz}


\subsection{Entradas}

Las entradas al reloj son 2:

\begin{enumerate}
    \item Tierra
    \item Voltaje de alimentación
\end{enumerate}

\subsection{Salidas}

Las salidas al reloj son 2: la señal de reloj, una onda cuadrada que alterna entre 0V y el voltaje de alimentación.

\begin{enumerate}
    \item Tierra
    \item Señal de reloj
\end{enumerate}

\newpage

\section{Secuenciador}

Para el módulo de secuenciador proponemos el uso del chip 4017, explicado en \ref{sec:chip4017}.


\begin{circuitikz}
    \draw (0,0) node[draw, minimum height=2.0cm, minimum width=1.0cm] (box) {Secuenciador};
    \draw (box.north) -- ++(0,0.2) node[anchor=south] {$V_{CC}$};
    \draw (box.south)   -- ++(0,-0.2) node[anchor=north] {Tierra};
    \draw (box.west) -| ++(-0.2,0) node[anchor=east] {Entrada de reloj};
    \draw (box.east) |- ++(0.4,+0.8) node[anchor=west] {Señal de paso 1};
    \draw (box.east) |- ++(0.4,+0.4) node[anchor=west] {Señal de paso 2};
    \draw (box.east) |- ++(0.4,+0.0) node[anchor=west] {Señal de paso 3};
    \draw (box.east) |- ++(0.4,-0.4) node[anchor=west] {Señal de paso 4};
    \draw (box.east) |- ++(0.4,-0.8) node[anchor=west] {Tierra};
\end{circuitikz}


\subsection{Entradas}

Las entradas al secuenciador son 3

\begin{enumerate}
    \item Tierra
    \item Voltaje de alimentación
    \item Señal de reloj
\end{enumerate}

\subsection{Salidas}

Las salidas al secuenciador son 5: tierra, y 1 por cada paso del secuenciador, alternando entre 0V y el voltaje de alimentación.

\begin{enumerate}
    \item Tierra
    \item Paso 1
    \item Paso 2
    \item Paso 3
    \item Paso 4
\end{enumerate}

\newpage

\section{Osciladores}

Para el módulo de osciladores proponemos el uso del chip 4093, explicado en \ref{sec:chip4093}.

\begin{circuitikz}
    \draw (0,0) node[draw, minimum height=2.0cm, minimum width=1.0cm] (box) {Osciladores};
    \draw (box.north) -- ++(0,0.2) node[anchor=south] {$V_{CC}$};
    \draw (box.south)   -- ++(0,-0.2) node[anchor=north] {Tierra};
    \draw (box.east) |- ++(0.4,+0.8) node[anchor=west] {Señal de oscilador 1};
    \draw (box.east) |- ++(0.4,+0.4) node[anchor=west] {Señal de oscilador 2};
    \draw (box.east) |- ++(0.4,+0.0) node[anchor=west] {Señal de oscilador 3};
    \draw (box.east) |- ++(0.4,-0.4) node[anchor=west] {Señal de oscilador 4};
    \draw (box.east) |- ++(0.4,-0.8) node[anchor=west] {Tierra};
\end{circuitikz}

\subsection{Entradas}


\begin{enumerate}
    \item Tierra
    \item Voltaje de alimentación
    \item Señal de paso 1
    \item Señal de paso 2
    \item Señal de paso 3
    \item Señal de paso 4
\end{enumerate}

\subsection{Salidas}

\begin{enumerate}
    \item Tierra
    \item Señal de oscilador 1
    \item Señal de oscilador 2
    \item Señal de oscilador 3
    \item Señal de oscilador 4
\end{enumerate}

\newpage

\section{Mezclador y amplificador}

Para el módulo de mezclador y amplificador proponemos el uso del chip 386.

\subsection{Entradas}

\begin{enumerate}
    \item Tierra
    \item Voltaje de alimentación
    \item Señal de audio 1
    \item Señal de audio 2
    \item Señal de audio 3
    \item Señal de audio 4
\end{enumerate}

\subsection{Salidas}

\begin{enumerate}
    \item Tierra
    \item Señal de audio mono amplificada
\end{enumerate}

\newpage

\section{Parlante}

Para parlante proponemos el uso de un parlante de 8 $\Omega$.

\begin{circuitikz}
    \draw (0,0) node[draw, minimum height=2.0cm, minimum width=1.0cm] (box) {Parlante};
    \draw (box.north) -- ++(0,0.2) node[anchor=south] {$V_{CC}$};
    \draw (box.south)   -- ++(0,-0.2) node[anchor=north] {Tierra};
    \draw (box.west) -| ++(-0.2,0) node[anchor=east] {Entrada de audio};
    \draw (box.east) |- ++(0.4,0) node[anchor=west] {Sonido};
\end{circuitikz}

\subsection{Entradas}

\begin{enumerate}
    \item Tierra
    \item Señal de audio monófonica
\end{enumerate}

\subsection{Salidas}

\begin{enumerate}
    \item Sonido
\end{enumerate}

\newpage

\end{document}
